\clearpage
\fancyhead[L]{\small PASQUICH. VINDICATION.}
\hypertarget{native-639-lower}{}
\expandafter\def\csname HistoricalNoteText35\endcsname{The collimation error of the Königsberg meridian circle amounts to nearly 2 degrees. Instrument makers usually fit their instruments with devices for eliminating collimation errors. The astronomer does not thank them for it. The fact that the reticle of some meridian circles can also be moved vertically should not be regarded as a virtue.\hfill G.}
\expandafter\def\csname HistoricalNoteText36\endcsname{For the first observation I relied on the hour angle, since the right ascension is evidently marred by a printing or writing error, and where it reads $357^\circ\,49'\,14''5$ it should read $357^\circ\,49'\,4''5$.}
\begin{center}\small Astronomische Nachrichten. Vol. III. No. 53. Pp. 77..84. 3 March 1824.\\[3pt]\rule{.72\linewidth}{.3pt}\end{center}
\begin{center}\emph{Vindication.}\end{center}

In one of the latest issues of the Correspondance Astronomique, which appears in Genoa, Professor \textsc{Pasquich} is accused of having fabricated the positions of the comet of 1821 printed in No. 2 of the Astr. Nachr., that is, of having calculated them from elements instead of deriving them from observations with the equatorial instrument, as \textsc{Pasquich} himself has stated he did. In support of this charge, the accuser does not claim that on the days for which the comet's positions were published, observations at the Ofen Observatory were either not made at all or made incorrectly. On the contrary, the accuser himself took part in the observations and himself produces some of the original observations from \textsc{Pasquich}'s own manuscript, authenticated by a notary. Nor does the accuser prove, or, strictly speaking, even claim, that the original observations, if correctly and fully reduced, would yield positions different from those published by \textsc{Pasquich}. On what grounds, then, one may ask, is an accusation founded in which a subordinate portrays his superior as guilty of the most dishonourable crime that can be laid to an astronomer's charge? Solely on the claim that the instrument had not been adjusted at all, that correct positions cannot be derived from observations with an unadjusted instrument, and that the positions published by \textsc{Pasquich}, whose accuracy Dr. \textsc{Ursin}'s calculation had immediately confirmed, must therefore have been dishonestly fabricated.

Astronomers need only this simple account, or a moderately attentive reading of the charge itself, to see that it is entirely without foundation. Every astronomer knows that positions derived from observations with an unadjusted instrument can be just as reliable as those derived from observations with a perfectly adjusted one. One determines by suitable observations the extent to which the instrument deviates from perfect adjustment and applies the resulting corrections to all other observations. Calculation is always the most precise method of correction. Whether the \SourceLeaf{640} corrections amount to a few seconds or a few minutes makes no difference to the accuracy of the results. It is for convenience, rather than for accuracy, that one generally prefers to deal only with small corrections, and for deviations that always have the same effect, even this scarcely matters at all\HistoricalNote{35}.

The accuser's assertion that correct positions cannot be derived from observations with an unadjusted instrument thus amounts only to an admission that needs no comment. His superior, who has long enjoyed a good reputation as a mathematician, could naturally have had no difficulty calculating corrections that are within the ability of even a beginner with a little practice.

On a cursory reading of the charge, one might perhaps think that the instrument had been very far from perfectly adjusted. Yet

\emph{First}, this does not follow from the alleged evidence. It is reported that once at twilight \textsc{Pasquich} was unable to find the comet with the equatorial instrument when the accuser could already see it in the comet seeker. This proves nothing at all. The comet could indeed have been in the field and yet, precisely because of the twilight, could still have been indiscernible at the high magnification when a comet seeker with good light-gathering power and lower magnification was already showing it. Likewise, the accuser's calculations on the observations do nothing in themselves to establish that the adjustment was deficient. In fact, with \emph{such} calculations he would have had to find a deviation even if none existed at all. For he neglects refraction altogether and, in his own fashion, makes fun of \textsc{Pasquich}'s assurance that he took it into account in reducing the observations. The accuser claims that refraction cannot, at its maximum, amount to as much as two seconds either in the difference of the right ascensions or in that of the declinations. Did the accuser not know how to perform even this easy calculation? On 27 February, the difference amounts to more than a minute in both cases.

\emph{Secondly}, however, the observations presented do in fact show clearly enough that the deviation of the instrument was by no means so enormously large.

The few observations of $\gamma$ Pegasi that the accuser has brought to our attention are indeed insufficient for a precise and complete determination of the correction elements of the instrument; they do suffice, however, to prove that the deviations are not so large as to make it difficult to locate stars; they suffice, moreover, to reduce the comet observations themselves to a very close approximation. I give the results of my calculation on this matter with all the greater pleasure, since they show with all the clarity one could wish that the accusation is more than groundless, that it is false.

The analysis of all 5 observations of $\gamma$ Pegasi gave me
\begin{center}\footnotesize\setlength{\tabcolsep}{3pt}
\begin{tabular}{@{}l r@{}}
Distance of the instrument pole from the true celestial pole & $=5'\;5''56$\\
Hour angle of the former . . . . . . . . . . . & $212^\circ\;27'\;27''$\\
Index correction for the hour angle . . . . . & $3'\;41''57$\\
for the declination . . . . . . . . . . . . . . & $4\;57.55$
\end{tabular}\end{center}

Whether and by how much the optical axis departs from parallelism with the declination circle, and the latter from perpendicularity to the equatorial circle, cannot be \SourceLeaf{641} determined from observations of a single star; but given the small difference between the declinations, this can have no perceptible effect on the reduction of the comet's positions.

The 5 observations of $\gamma$ Pegasi are represented by these elements
as follows:

\begin{center}\footnotesize\setlength{\tabcolsep}{4pt}
\begin{tabular}{r|r|r|l|r}
 & \multicolumn{2}{c|}{Corr. to hr. ang. for} & Corr. observed & Diff. from\\
 & instr. & refr. & hour angle & true value \\
\hline
Feb. 20 & $-162''99$ & $+147''26$ & $81^\circ\,36'\,45''01$ & $+5''42$ \\
22 & $-174.81$ & $+95.54$ & $69\ \ 38\ \ 14.47$ & $+3.23$ \\
22 & $-168.76$ & $+116.59$ & $75\ \ 29\ \ 8.57$ & $+0.33$ \\
27 & $-174.97$ & $+95.56$ & $69\ \ 28\ \ 41.31$ & $+1.42$ \\
27 & $-168.59$ & $+117.84$ & $75\ \ 39\ \ 9.99$ & $-10.40$ \\
\end{tabular}
\end{center}

\begin{center}\footnotesize\setlength{\tabcolsep}{4pt}
\begin{tabular}{r|r|r|l|r}
 & \multicolumn{2}{c|}{Corr. to decl. for} & Corr. observed & Diff. from\\
 & instr. & refr. & declination & true value \\
\hline
Feb. 20 & $-97''54$ & $-147''48$ & $14^\circ\,11'\,19''98$ & $-0''63$ \\
22 & $-54.04$ & $-96.10$ & $14\ \ 11\ \ 21.86$ & $+1.42$ \\
22 & $-74.06$ & $-116.33$ & $14\ \ 11\ \ 17.61$ & $-2.83$ \\
27 & $-53.50$ & $-96.15$ & $14\ \ 11\ \ 20.35$ & $+0.33$ \\
27 & $-74.68$ & $-117.59$ & $14\ \ 11\ \ 21.73$ & $+1.71$ \\
\end{tabular}
\end{center}

Applying the same elements to the reduction of the comet observations
gives

\begin{center}\footnotesize\setlength{\tabcolsep}{4pt}
\emph{For the hour angle}
\end{center}

\begin{center}\footnotesize\setlength{\tabcolsep}{4pt}
\begin{tabular}{r|r|r|l|r}
 & \multicolumn{2}{c|}{Correction for} & Corrected & \shortstack{Dev. of means\\from \textsc{Pasquich}'s}\\
 & instr. & refr. & hour angle & values \\
\hline
Feb. 22 & $-171''15$ & $+101''34$ & $71^\circ\,48'\,50''19$ & \raisebox{-5.5pt}[0pt][0pt]{$-4''34$} \\
22 & $-165.09$ & $+124.21$ & $77\ \ 39\ \ 19.12$ & \\
26 & $-164.37$ & $+134.96$ & $79\ \ 28\ \ 15.59$ & \raisebox{-5.5pt}[0pt][0pt]{$+1.37$} \\
26 & $-161.37$ & $+154.53$ & $82\ \ 44\ \ 53.16$ & \\
27 & $-161.34$ & $+157.80$ & $83\ \ \ 9\ \ 56.48$ & \raisebox{-5.5pt}[0pt][0pt]{$+4.95$} \\
27 & $-157.98$ & $+191.42$ & $87\ \ 10\ \ 33.44$ & \\
\end{tabular}
\end{center}

\begin{center}\footnotesize\setlength{\tabcolsep}{4pt}
\emph{For the declination}
\end{center}

\begin{center}\footnotesize\setlength{\tabcolsep}{4pt}
\begin{tabular}{r|r|r|l|r}
 & \multicolumn{2}{c|}{Correction for} & Corrected & \shortstack{Dev. of means\\from \textsc{Pasquich}'s}\\
 & instr. & refr. & declination & values \\
\hline
Feb. 22 & $-61''15$ & $-100''78$ & $14^\circ\,38'\,10''07$ & \raisebox{-5.5pt}[0pt][0pt]{$+7''42$} \\
22 & $-82.05$ & $-123.18$ & $14\ \ 38\ \ 24.77$ & \\
26 & $-89.01$ & $-134.43$ & $14\ \ 23\ \ \ 9.56$ & \raisebox{-5.5pt}[0pt][0pt]{$+0.91$} \\
26 & $-102.09$ & $-154.65$ & $14\ \ 23\ \ 18.26$ & \\
27 & $-103.80$ & $-158.18$ & $14\ \ 18\ \ 44.02$ & \raisebox{-5.5pt}[0pt][0pt]{$-5.58$} \\
27 & $-120.75$ & $-193.94$ & $14\ \ 18\ \ 34.82$ & \\
\end{tabular}
\end{center}

For anyone who is no stranger to astronomical calculation, these deviations
must be the clearest proof that \textsc{Pasquich}'s positions really were derived
from the observations by proper reduction calculations. To be sure, as shown above, the
whole accusation gives no grounds for doubting this. To be sure, even quite
apart from this accusation, there is no ground for such a doubt, which would
be absurd in itself if only because it takes much less work to reduce the
actual observations than to calculate the positions of the comet from
elements. Yet in this case the observation itself really does provide the
most evident proof of the genuine origin of the comet positions published
by \textsc{Pasquich}. For most readers the smallness of the deviations found
above will itself be such proof, but much \SourceLeaf{642}\enlargethispage{1pt} stronger proof lies in their \emph{regular pattern}.
Our reduction elements may differ considerably from the true ones. This
must have some effect on the comet positions, though always only a small one.
Comparing the results of calculations made with two different systems of
reduction elements applied to one and the same series of observations will
therefore show small differences, but differences that necessarily follow
a regular pattern. Comparing positions that had been calculated from
elements with positions derived from actual observations, on the other
hand, would yield differences that would inevitably bear the unmistakable
trace, in absolute terms, of the errors in the elements and, in relative
terms, of the irregular errors of observation (together with any small
changes deliberately and arbitrarily made, if one were to suppose there
were any). It would be more than a miracle if someone so unreasonable as
to fabricate positions from elements when he has good observations before him that
are easy to reduce were, in the course of such a foolish fraud, to draw
such a winning five-number combination from the lottery urn that he hit
on exactly what reducing his observations would have given him. As for
the above differences, the extent of the regularity in their pattern is
best seen by deriving the reduction elements from the observations of the comet
itself, compared with the published positions, instead of from the observations
of $\gamma$ Pegasi, and by using only two observations of the comet and then
seeing how well the third agrees with them. In this way one finds\HistoricalNote{36}:

\begin{center}\footnotesize\setlength{\tabcolsep}{3pt}
\begin{tabular}{@{}l r@{}}
Distance of the instrument's pole from the true celestial pole & $2'\;32''41$\\
Hour angle of the former . . . . . . . . . . & $170^\circ\;25\;17''$\\
Index correction for the hour angle . . . & $-3\;23.47$\\
Index correction for the declination . . . . . & $-1\;34.38$
\end{tabular}\end{center}

The reduction of the observations of the comet, carried out in a manner
similar to that above, then gives

\begin{center}\footnotesize\setlength{\tabcolsep}{3pt}\begin{tabular}{r | r r | r | r r | r}
 & \multicolumn{2}{c|}{Hour angle} & \shortstack{according to\\\textsc{Pasquich}} & \multicolumn{2}{c|}{Declination} & \shortstack{according to\\\textsc{Pasquich}} \\
\hline
Feb. 22 & $74^\circ\,44'$ & $8''99$ & $74^\circ\,44'\,9''$ & $14^\circ\,38'$ & $9''82$ & $14^\circ\,38'\,10''$ \\
26 & $81$ & $6\ \ 32.98$ & $81\ \ 6\ \ 33$ & $14$ & $23\ \ 13.36$ & $14\ \ 23\ \ 13$ \\
27 & $85$ & $10\ \ 10.03$ & $85\ \ 10\ \ 10$ & $14$ & $18\ \ 44.81$ & $14\ \ 18\ \ 45$ \\
\end{tabular}\end{center}

thus giving truly \emph{perfect} agreement, since \textsc{Pasquich} did not
specify fractions of seconds.

This system of reduction elements represents the observations of $\gamma$ Pegasi
closely enough, even if not quite as closely as the one above. The differences are

\begin{center}\footnotesize\setlength{\tabcolsep}{4pt}\begin{tabular}{r | r | r}
 & \shortstack{in the\\hour angle} & \shortstack{in the\\declination} \\
\hline
Feb. 20 & $+3''51$ & $-0''39$ \\
22 & $+12.47$ & $-10.32$ \\
22 & $+4.08$ & $-9.82$ \\
27 & $+10.81$ & $-11.42$ \\
27 & $-6.80$ & $-5.10$ \\
\end{tabular}\end{center}

We do not know the reduction elements that \textsc{Pasquich} actually used.
They may differ considerably from those just given, for it is too uncertain
a matter to infer from the corrections \SourceLeaf{643} the elements on which
they are based, especially since the published data do not allow us to obtain
the corrections for the instrument separately from the refraction, and
\textsc{Pasquich} presumably calculated the refractions from other tables
and perhaps by other methods. Nor is the question whether the reduction elements
that \textsc{Pasquich} used were the most accurate possible at all relevant to
the matter. Or rather, merely to raise this question is already to acquit
\textsc{Pasquich} of the charge brought against him. Still, I am glad to declare
that I see no reason at all to doubt that \textsc{Pasquich} derived his reduction
elements with due deliberation from observations that were probably far more
numerous and perhaps entirely different, since the above calculation itself
makes it perfectly clear that \textsc{Pasquich} really did reduce his observations
as absolute and not as differential observations.

So far I have considered the matter only from the scientific point of view,
but how could one fail to consider it also from the standpoint of honour and
integrity and to share fully the just indignation that the noble \textsc{Olbers}
expressed so aptly!

\noindent Göttingen 1824 March 3. \hfill C.\,F.\ \textsc{Gauss}.

[Handwritten remark: \textsc{Daniel Kmeth}, b. in Bries (Brezno Bamja) in Hungary 1783 Jan. 15, died in Kaschau 1825 June 20.]
